\documentclass[12pt,a4paper]{article}
\usepackage[utf8]{inputenc}
\usepackage[T1]{fontenc}
\usepackage[spanish]{babel}
\usepackage{geometry}
\usepackage{graphicx}
\usepackage{float}
\usepackage{hyperref}
\usepackage{longtable}
\usepackage{array}
\usepackage{xcolor}
\usepackage{listings}
\usepackage{enumitem}

\geometry{margin=2.5cm}
\setlength{\parskip}{0.7em}
\setlength{\parindent}{0pt}
\hypersetup{colorlinks=true, linkcolor=blue, urlcolor=blue}

\lstdefinestyle{code}{
  basicstyle=\ttfamily\small,
  breaklines=true,
  frame=single,
  backgroundcolor=\color{gray!8}
}

\newcommand{\evidenceimage}[3]{
  \begin{figure}[H]
    \centering
    \IfFileExists{figures/#1}{\includegraphics[width=0.92\textwidth]{figures/#1}}{\fbox{\parbox[c][5cm][c]{0.86\textwidth}{\centering Captura pendiente: \texttt{figures/#1}}}}
    \caption{#2}
    \label{#3}
  \end{figure}
}

\begin{document}

\begin{titlepage}
  \centering
  {\Large Universidad Politecnica Salesiana\par}
  {\large Maestria en Ingenieria de Software\par}
  {\large Desarrollo de Aplicaciones Empresariales\par}
  \vfill
  {\Huge\bfseries Sistema de Reserva de Canchas Deportivas\par}
  \vspace{0.5cm}
  {\Large Informe academico final\par}
  \vfill
  \begin{tabular}{ll}
    Frontend 1: & Enzo Aliatis \\
    Frontend 2: & Jordan Murillo \\
    Backend 1: & Galo Suarez \\
    Backend 2: & Mauro Cadme \\
  \end{tabular}
  \vfill
  {\large Agosto 2026\par}
\end{titlepage}

\tableofcontents
\newpage

\section{Resumen ejecutivo}

El proyecto implementa un sistema web para reserva de canchas deportivas de padel, tenis y basquet. La solucion permite a usuarios finales registrarse, iniciar sesion, consultar disponibilidad, crear reservas, revisar su historial y cancelar reservas. Tambien permite a administradores gestionar canchas, usuarios, reservas globales, bloqueos de mantenimiento y reportes basicos de ocupacion.

La arquitectura se construyo con microfrontends Angular integrados mediante Module Federation, microservicios Spring Boot y una base PostgreSQL independiente por servicio. El sistema completo puede levantarse localmente con Docker Compose y expone documentacion Swagger/OpenAPI en los cuatro microservicios. La validacion funcional manual fue realizada; quedan pendientes las pruebas automatizadas end-to-end y la consolidacion formal de evidencias de prueba.

\section{Introduccion}

El proyecto corresponde al trabajo integrador de la asignatura Desarrollo de Aplicaciones Empresariales. Su proposito es aplicar conceptos de arquitectura empresarial moderna mediante una solucion acotada, verificable y desplegable localmente. El dominio elegido, reserva de canchas deportivas, permite trabajar reglas de negocio relevantes como disponibilidad, no solapamiento de horarios, roles diferenciados, estados de reserva y reportes operativos.

El alcance tecnico prioriza independencia entre componentes. En frontend se separan responsabilidades en un shell y tres remotes. En backend se dividen dominios en cuatro microservicios. En datos se mantiene la propiedad por servicio, evitando claves foraneas y consultas cruzadas entre bases de microservicios distintos.

\section{Objetivos}

\subsection{Objetivo general}

Desarrollar un sistema web de reserva de canchas deportivas que permita a usuarios finales y administradores operar los flujos principales del dominio, utilizando microfrontends con Module Federation, microservicios Spring Boot y persistencia PostgreSQL.

\subsection{Objetivos especificos}

\begin{itemize}
  \item Implementar un shell Angular que orqueste microfrontends remotos.
  \item Construir microservicios independientes para usuarios, canchas, reservas y reportes.
  \item Modelar una base PostgreSQL independiente por microservicio.
  \item Implementar autenticacion JWT y control de acceso por roles.
  \item Aplicar reglas de negocio de reserva, disponibilidad y cancelacion.
  \item Proveer despliegue local reproducible mediante Docker Compose.
  \item Documentar arquitectura, modelo de datos, APIs, pruebas y pendientes.
\end{itemize}

\section{Alcance funcional}

El sistema cubre los siguientes procesos principales:

\begin{itemize}
  \item Registro e inicio de sesion de usuarios.
  \item Consulta de canchas activas y horarios disponibles.
  \item Creacion de reservas por parte de usuarios finales.
  \item Consulta y cancelacion de reservas propias.
  \item Gestion administrativa de canchas, horarios, estados y bloqueos.
  \item Gestion administrativa de usuarios.
  \item Consulta administrativa de reservas globales.
  \item Reportes basicos de ocupacion, demanda y cancelaciones.
\end{itemize}

Quedan fuera de alcance pagos, notificaciones automaticas, reservas recurrentes, torneos, aplicacion movil nativa y reportes avanzados de inteligencia de negocio.

\section{Roles, permisos y modulos}

\subsection{Roles}

\begin{longtable}{p{3.5cm}p{10cm}}
\textbf{Rol} & \textbf{Permisos principales} \\
\hline
Usuario final & Registrarse, iniciar sesion, consultar disponibilidad, crear reservas, consultar sus reservas y cancelar reservas propias antes del inicio. \\
Administrador & Gestionar canchas, bloqueos, usuarios, reservas globales y reportes. Puede cancelar cualquier reserva del sistema. \\
\end{longtable}

\subsection{Modulos funcionales}

\begin{longtable}{p{4cm}p{4cm}p{6cm}}
\textbf{Modulo} & \textbf{Componente} & \textbf{Responsabilidad} \\
\hline
Shell & \texttt{frontend/shell} & Layout, autenticacion, navegacion, guards y carga de remotes. \\
Cliente & \texttt{mf-clientes} & Disponibilidad, reserva y consulta de reservas propias. \\
Administracion & \texttt{mf-administracion} & Canchas, usuarios y reservas administrativas. \\
Reportes & \texttt{mf-reportes} & Consulta de indicadores basicos para administradores. \\
\end{longtable}

\evidenceimage{01-landing.png}{Pantalla inicial del shell del sistema.}{fig:landing}
\evidenceimage{02-login.png}{Pantalla de inicio de sesion.}{fig:login}

\section{Reglas de negocio}

Las reglas de negocio implementadas se concentran en el proceso de reserva y en la administracion de canchas:

\begin{itemize}
  \item Una reserva debe asociarse a un usuario autenticado y a una cancha existente.
  \item No se permite crear reservas sobre horarios ocupados.
  \item La hora de inicio debe pertenecer al horario de atencion de la cancha.
  \item No se permite reservar una cancha inactiva.
  \item Los bloqueos de mantenimiento impiden crear reservas en el periodo bloqueado.
  \item Existe un limite configurable de reservas activas simultaneas por usuario.
  \item Solo el administrador puede crear, editar o inactivar canchas.
  \item Toda reserva mantiene un estado: \texttt{CONFIRMADA}, \texttt{CANCELADA} o \texttt{FINALIZADA}.
\end{itemize}

La regla de no solapamiento se refuerza en PostgreSQL mediante un constraint \texttt{EXCLUDE USING gist} sobre la tabla \texttt{reservas}. Esto evita condiciones de carrera que una validacion solamente en Java no podria prevenir de forma completa.

\section{Arquitectura implementada}

La arquitectura implementada combina microfrontends, microservicios y bases de datos independientes. El navegador carga el shell en el puerto 4300; este shell integra remotes en los puertos 4201, 4202 y 4203. Los microfrontends consumen directamente las APIs publicas de los microservicios. No se utiliza API Gateway ni BFF.

\evidenceimage{arquitectura-contenedores.png}{Vista de contenedores del sistema.}{fig:contenedores}

\begin{longtable}{p{4cm}p{4cm}p{6cm}}
\textbf{Capa} & \textbf{Tecnologia} & \textbf{Detalle} \\
\hline
Frontend & Angular 20, Module Federation & Shell y tres remotes. \\
Backend & Java 17, Spring Boot, Maven & Cuatro microservicios independientes. \\
Datos & PostgreSQL, Flyway & Una base de datos por servicio. \\
Despliegue local & Docker Compose & Levanta frontend, backend y bases de datos. \\
Documentacion API & Swagger/OpenAPI & Disponible en runtime por microservicio. \\
\end{longtable}

\section{Microfrontends}

El frontend se estructura en cuatro aplicaciones Angular:

\begin{itemize}
  \item \texttt{shell}: host de Module Federation, responsable del layout, autenticacion, rutas protegidas y navegacion.
  \item \texttt{mf-clientes}: remote de usuario final, enfocado en disponibilidad y reservas.
  \item \texttt{mf-administracion}: remote administrativo para canchas, usuarios y reservas.
  \item \texttt{mf-reportes}: remote administrativo para reportes.
\end{itemize}

Cada remote expone sus rutas mediante \texttt{./Routes}. El shell decide que remote cargar segun la ruta y el rol del usuario. Esta decision reduce acoplamiento entre modulos y permite construir y desplegar cada microfrontend de forma independiente.

\evidenceimage{03-cliente-disponibilidad.png}{Modulo de cliente para disponibilidad y reserva.}{fig:cliente-disponibilidad}
\evidenceimage{04-cliente-mis-reservas.png}{Modulo de cliente para consulta de reservas propias.}{fig:cliente-reservas}
\evidenceimage{05-admin-canchas.png}{Modulo administrativo de canchas.}{fig:admin-canchas}
\evidenceimage{06-admin-usuarios.png}{Modulo administrativo de usuarios.}{fig:admin-usuarios}
\evidenceimage{07-admin-reservas.png}{Modulo administrativo de reservas.}{fig:admin-reservas}
\evidenceimage{08-reportes.png}{Modulo administrativo de reportes.}{fig:reportes}

\section{Microservicios}

El backend esta compuesto por cuatro microservicios Spring Boot:

\begin{longtable}{p{3.5cm}p{2cm}p{3cm}p{6cm}}
\textbf{Servicio} & \textbf{Puerto} & \textbf{Base} & \textbf{Responsabilidad} \\
\hline
\texttt{ms-usuarios} & 8081 & \texttt{usuarios\_db} & Registro, login JWT, usuarios y roles. \\
\texttt{ms-canchas} & 8082 & \texttt{canchas\_db} & Canchas, horarios, deportes y mantenimientos. \\
\texttt{ms-reservas} & 8083 & \texttt{reservas\_db} & Disponibilidad, reserva, cancelacion y estados. \\
\texttt{ms-reportes} & 8084 & \texttt{reportes\_db} & Reportes basicos y auditoria de consultas. \\
\end{longtable}

\subsection{Endpoints principales}

\begin{longtable}{p{4cm}p{10cm}}
\textbf{Servicio} & \textbf{Endpoints publicos principales} \\
\hline
\texttt{ms-usuarios} & \texttt{POST /api/v1/auth/registro}, \texttt{POST /api/v1/auth/login}, \texttt{GET /api/v1/usuarios/me}, \texttt{GET /api/v1/usuarios}, \texttt{PATCH /api/v1/usuarios/\{id\}/estado}. \\
\texttt{ms-canchas} & \texttt{GET /api/v1/canchas}, \texttt{POST /api/v1/canchas}, \texttt{PUT /api/v1/canchas/\{id\}}, \texttt{PATCH /api/v1/canchas/\{id\}/estado}, \texttt{POST /api/v1/canchas/\{id\}/mantenimientos}. \\
\texttt{ms-reservas} & \texttt{GET /api/v1/disponibilidad}, \texttt{POST /api/v1/reservas}, \texttt{GET /api/v1/reservas/mias}, \texttt{GET /api/v1/reservas}, \texttt{PATCH /api/v1/reservas/\{id\}/cancelacion}. \\
\texttt{ms-reportes} & \texttt{GET /api/v1/reportes/ocupacion}. \\
\end{longtable}

\evidenceimage{endpoints.png}{Resumen visual de endpoints principales.}{fig:endpoints}

\section{Modelo de datos}

Cada microservicio posee su propia base PostgreSQL. No existen joins, claves foraneas ni consultas entre bases de microservicios distintos. Las relaciones externas se guardan como UUID, por ejemplo \texttt{usuario\_id} y \texttt{cancha\_id} en \texttt{reservas\_db}.

\evidenceimage{modelo-datos.png}{Modelo de datos por servicio.}{fig:modelo-datos}

\subsection{usuarios\_db}

\begin{lstlisting}[style=code]
usuarios(id, nombre, email, password_hash, rol, activo, creado_en, actualizado_en)
\end{lstlisting}

\subsection{canchas\_db}

\begin{lstlisting}[style=code]
canchas(id, nombre, deporte, hora_apertura, hora_cierre, activa, creado_en, actualizado_en)
bloqueos_mantenimiento(id, cancha_id, inicio, fin, motivo, creado_en)
\end{lstlisting}

\subsection{reservas\_db}

\begin{lstlisting}[style=code]
reservas(id, usuario_id, cancha_id, cancha_nombre, deporte, inicio, fin, estado,
         cancelada_por, motivo_cancelacion, creado_en, actualizado_en)
\end{lstlisting}

\subsection{reportes\_db}

\begin{lstlisting}[style=code]
reporte_consultas(id, usuario_id, fecha_desde, fecha_hasta, total_reservas, generado_en)
\end{lstlisting}

\section{Infraestructura y despliegue local}

El despliegue local se realiza con \texttt{infra/docker-compose.yml}. El archivo levanta cuatro contenedores PostgreSQL, cuatro microservicios, los tres microfrontends remotos y el shell.

\begin{lstlisting}[style=code]
cp infra/.env.example infra/.env
docker compose -f infra/docker-compose.yml up --build
\end{lstlisting}

La aplicacion se utiliza desde \url{http://localhost:4300}. Los puertos 4201, 4202 y 4203 corresponden a los remotes y se usan principalmente para inspeccion aislada.

\evidenceimage{11-docker-compose.png}{Evidencia de despliegue local con Docker Compose.}{fig:compose}

\section{Seguridad e integracion}

La autenticacion de usuarios se realiza con JWT emitido por \texttt{ms-usuarios}. El shell almacena la sesion localmente y usa guards para controlar la navegacion por rol. La autorizacion efectiva se valida en backend mediante configuracion de seguridad Spring.

La comunicacion interna entre microservicios usa endpoints bajo \texttt{/internal/**} protegidos con la cabecera \texttt{X-Service-Key}. Esta llave separa el trafico interno del trafico de usuarios finales.

Las credenciales incluidas en archivos de ejemplo son solo para demostracion local. En ambientes reales deben reemplazarse por variables de entorno no versionadas.

\section{Pruebas y validacion}

El proyecto incluye pruebas unitarias iniciales en backend y frontend, Swagger/OpenAPI disponible en runtime y una coleccion Postman con endpoints y flujo funcional completo.

\begin{longtable}{p{4cm}p{6cm}p{4cm}}
\textbf{Tipo} & \textbf{Evidencia} & \textbf{Estado} \\
\hline
Unitarias backend & Tests en \texttt{backend/ms-*/src/test} & Implementadas inicialmente. \\
Unitarias frontend & Specs Angular en shell y remotes & Implementadas inicialmente. \\
API manual & Coleccion Postman & Versionada. \\
Swagger/OpenAPI & \texttt{/swagger-ui.html} en cada servicio & Disponible en runtime. \\
Funcional manual & Flujos principales levantados con Docker Compose & Validado manualmente. \\
End-to-end automatizado & Cypress, Playwright u otra herramienta & Pendiente. \\
\end{longtable}

\evidenceimage{09-swagger-usuarios.png}{Swagger UI de ms-usuarios.}{fig:swagger-usuarios}
\evidenceimage{10-swagger-reservas.png}{Swagger UI de ms-reservas.}{fig:swagger-reservas}
\evidenceimage{12-postman-collection.png}{Coleccion Postman del backend.}{fig:postman}

\section{Entregables}

\begin{longtable}{p{1.5cm}p{9cm}p{3cm}}
\textbf{Codigo} & \textbf{Descripcion} & \textbf{Estado} \\
\hline
E1 & Documento de arquitectura con microfrontends, microservicios y modelo de datos. & En preparacion en LaTeX. \\
E2 & Shell y microfrontends funcionando integrados por Module Federation. & Implementado. \\
E3 & Microservicios Spring Boot documentados con Swagger/OpenAPI y coleccion de pruebas. & Implementado. \\
E4 & Scripts DDL y seed de PostgreSQL. & Implementado. \\
E5 & Manual breve de despliegue local. & En preparacion en LaTeX. \\
E6 & Presentacion final y demostracion en vivo. & Pendiente. \\
\end{longtable}

\section{Estado actual y pendientes}

El sistema esta implementado, integrado y levantado correctamente en entorno local. La validacion funcional manual esta realizada. Los pendientes principales son:

\begin{itemize}
  \item Crear y ejecutar pruebas automatizadas end-to-end.
  \item Consolidar evidencias formales de pruebas.
  \item Exportar o adjuntar contratos OpenAPI si se requiere entrega estatica.
  \item Preparar presentacion final y guion de demostracion.
  \item Renderizar los documentos PDF desde estas fuentes LaTeX.
\end{itemize}

\section{Conclusiones}

El proyecto cumple el objetivo principal de implementar una aplicacion empresarial distribuida con separacion clara entre frontend, backend y datos. La division en microfrontends permite aislar flujos de cliente, administracion y reportes, mientras que los microservicios separan los dominios de usuarios, canchas, reservas y reportes.

La independencia de datos se respeto mediante una base PostgreSQL por servicio. La regla critica de no solapamiento de reservas fue reforzada en base de datos con un constraint especializado, lo cual mejora la consistencia del sistema bajo concurrencia. La decision de no utilizar API Gateway mantiene la arquitectura simple y suficiente para el alcance academico definido.

\section{Anexos}

\subsection{Equipo}

\begin{longtable}{p{3cm}p{5cm}p{5cm}}
\textbf{Rol} & \textbf{Responsabilidad primaria} & \textbf{Integrante} \\
\hline
Frontend 1 & Shell y autenticacion en interfaz; apoyo a mf-clientes. & Enzo Aliatis \\
Frontend 2 & mf-administracion y mf-reportes. & Jordan Murillo \\
Backend 1 & ms-usuarios y ms-canchas. & Galo Suarez \\
Backend 2 & ms-reservas y ms-reportes. & Mauro Cadme \\
\end{longtable}

\subsection{URLs locales}

\begin{longtable}{p{5cm}p{8cm}}
\textbf{Componente} & \textbf{URL} \\
\hline
Shell & \url{http://localhost:4300} \\
mf-clientes & \url{http://localhost:4201} \\
mf-administracion & \url{http://localhost:4202} \\
mf-reportes & \url{http://localhost:4203} \\
Swagger ms-usuarios & \url{http://localhost:8081/swagger-ui.html} \\
Swagger ms-canchas & \url{http://localhost:8082/swagger-ui.html} \\
Swagger ms-reservas & \url{http://localhost:8083/swagger-ui.html} \\
Swagger ms-reportes & \url{http://localhost:8084/swagger-ui.html} \\
\end{longtable}

\subsection{Repositorio}

El codigo fuente se mantiene en un monorepo publico con carpetas principales \texttt{frontend/}, \texttt{backend/}, \texttt{database/}, \texttt{infra/} y \texttt{docs/}.

\end{document}

